\documentclass[11pt,a4paper]{article}
\usepackage[T1]{fontenc}
\usepackage[utf8]{inputenc}
\usepackage[main=italian,provide=*]{babel}
\usepackage{amsmath,amssymb}
\usepackage{geometry}
\geometry{margin=2.5cm}
\usepackage{booktabs}
\usepackage{seqsplit}
\usepackage{hyperref}
\hypersetup{colorlinks=true,linkcolor=blue,citecolor=blue,urlcolor=blue}

\title{VQE (con e senza termine DM) — dimero}
\author{Note di lavoro — Tesi triennale, Samuele Schiavi\\ Relatore: Prof.\ Alessandro Chiesa}
\date{}

\begin{document}
\maketitle

\section*{Introduzione}

Questo documento raccoglie i file della fase VQE per il dimero di spin-1/2 ($N=2$).

\textbf{Nota sullo scope, a differenza dei documenti gemelli di anello e catena.}
Per i due trimeri il termine DM è stato aggiunto in una fase successiva e ben
distinta (i moduli VQE senza DM hanno \texttt{dm\_mode}/\texttt{D} come parametri
opzionali con default che li escludono, e gli script principali usano solo $D=0$).
Per il dimero questa separazione \textbf{non esiste}: il DM è stato parte integrante
del lavoro fin dall'inizio, e l'output standard di \texttt{vqe\_dimer.py} è già il
confronto a quattro pannelli HA/PMA $\times$ $D{=}0$/$D{=}0.2$. Non c'è quindi un
sottoinsieme "senza DM" da isolare senza snaturare il contenuto — questo documento
copre perciò tutta la fase VQE del dimero in un unico file, non sdoppiato.

Per ciascun file: cosa fa, con quali verifiche è stato controllato in questa sessione,
e cosa se ne può concludere — per il modulo \texttt{.py}, anche le funzioni interne e
come si usa il file.

\section*{\texttt{vqe\_dimer.py} + \texttt{vqe\_dimer.png}}

Modulo di produzione principale: ansatz HA (hardware-efficient, 6 parametri) e PMA
(1 parametro), entrambi confrontati a $D=0$ e $D=0.2$ nello stesso sweep — quattro
righe della griglia risultante (HA $D{=}0$, HA $D{=}0.2$, PMA $D{=}0$, PMA $D{=}0.2$).

\textbf{Funzioni interne.}
\begin{itemize}
\item \texttt{make\_ha\_ansatz(reps)} — l'ansatz hardware-efficient, $2(\text{reps}{+}1)$
parametri.
\item \texttt{make\_pma\_ansatz(b,\ J,\ K)} — l'ansatz fisicamente motivato, blocco
$W_{01}(\theta)$ (mirror dell'ansatz "originale" di Crippa et al.), $K$ parametri.
\item \texttt{\_energy(params,\ ansatz,\ hamiltonian,\ estimator)} — la funzione
obiettivo per \texttt{\seqsplit{scipy.optimize.minimize}}.
\item \texttt{run\_vqe(b,\ J,\ D,\ ansatz\_type,\ reps,\ K,\ n\_restarts,\ method,\
seed)} — VQE a un singolo punto, con $D$ come parametro pieno (non opzionale
escluso): qui il DM è sempre potenzialmente attivo, non un'aggiunta successiva.
\item \texttt{vqe\_sweep(b\_values,\ J,\ D,\ ansatz\_type,\ reps,\ K,\ n\_restarts,\
method,\ seed,\ verbose)} — chiama \texttt{run\_vqe} su una griglia di campi $b$.
\item \texttt{\_plot\_row(axes,\ results,\ J,\ row\_label,\ color)} — disegna una
riga della griglia (energia, errore, fidelity, magnetizzazione).
\item \texttt{plot\_grid(all\_results,\ J,\ base\_path)} — la griglia completa
4 righe $\times$ 4 colonne, salvata in \texttt{.png}/\texttt{.pdf}.
\end{itemize}

\textbf{Come si usa.} Come libreria: \texttt{from vqe\_dimer import run\_vqe,
vqe\_sweep, \ldots} — è la dipendenza comune a \texttt{vqe\_dimer\_spiegato.ipynb} e
a \texttt{vqe\_ground\_state\_test2.ipynb}. Come script standalone: \texttt{python3
vqe\_dimer.py} esegue lo sweep completo per le quattro configurazioni e salva
\texttt{vqe\_dimer.png}/\texttt{.pdf}. Dipende da \texttt{dimer\_exact.py} nella
stessa cartella.

\textbf{Verifiche.} Rieseguito per intero in questa sessione: fidelity $\mathcal
F=1$ esatto per HA (entrambi $D$) e per PMA a $D=0$; per PMA a $D=0.2$ (1 solo
parametro) fidelity minima $\approx0.82$ vicino a $B/J=2$ — coerente con il crollo
già noto e documentato (conservazione di $M$ rotta solo parzialmente da un ansatz a
un parametro). Figura rigenerata e confrontata pixel per pixel con quella allegata:
identica.

\textbf{Conclusione.} Modulo maturo, nucleo di tutta la fase VQE del dimero.

\section*{\texttt{vqe\_dimer\_spiegato.ipynb}}

Guida interattiva al modulo precedente: principio variazionale, come l'estimator
calcola $\langle H\rangle_{\boldsymbol\theta}$, i due ansatz HA vs PMA (con
derivazione esplicita del blocco $W_{01}(\theta)$), perché il PMA a 1 parametro
funziona a $D=0$ ma non a $D=0.2$ (la stessa domanda che ha poi motivato
\texttt{analisi\_rotazioni\_PMA.ipynb}).

\textbf{Verifiche.} Eseguito per intero in questa sessione (31 celle, zero errori).
Il file allegato differisce dalla versione già nota solo per una cella vuota in più
in fondo — differenza non sostanziale.

\textbf{Conclusione.} Guida completa, già alla base della struttura poi replicata
per anello e catena (\texttt{vqe\_trimer\_ring\_spiegato.ipynb\ e affini}).

\section*{\texttt{confronto\_ansatz\_entangler.ipynb}}

Notebook di confronto sistematico: 12 ansatz (quattro generici A–D, famiglie
PMA-1q e PMA-2q a diversi conteggi di parametri), fidelity a $D=0$ e al crescere di
$D$. Stabilisce \texttt{PMA-2q.3} come ansatz canonico (minimo teorico di
parametri, zero ridondanza) — la stessa scelta poi ereditata da tutti i moduli
successivi del progetto, dimero incluso.

\textbf{Verifiche.} Il tentativo di riesecuzione completa in questa sessione è
stato interrotto (il processo di calcolo, pesante per il multistart su 12 ansatz,
non ha completato prima di un reset dell'ambiente). Il file conserva però gli
output di un'esecuzione precedente, completa e senza errori: controllati
puntualmente più valori (tabella parametri/ansatz, raccomandazione
\texttt{PMA-2q.3}, il caso limite dei layer RBS $M$-conservanti che restano bloccati
a $\mathcal F=0.5082$ qualunque sia $K$) — tutti coerenti con quanto già
documentato altrove nel progetto.

\textbf{Considerazioni.} A differenza degli altri file di questa sessione, per
questo non è stato possibile ottenere una riesecuzione fresca end-to-end — la
verifica qui si basa sulla coerenza interna e incrociata degli output già presenti,
non su un'esecuzione osservata direttamente in questa sessione.

\textbf{Conclusione.} È il notebook da cui discende la scelta dell'ansatz canonico
usata ovunque nel progetto — vale la pena, in futuro, una riesecuzione completa
dedicata per avere anche qui una verifica diretta, non solo di coerenza.

\section*{\texttt{analisi\_rotazioni\_PMA.ipynb}}

Indagine mirata: dove mettere le rotazioni che rompono la conservazione di $M$ (un
qubit o due? parametro condiviso o indipendente?), con derivazioni ed esperimenti
che falliscono apposta per isolare la causa. Risolve la domanda rimasta aperta da
\texttt{vqe\_dimer\_spiegato.ipynb} (perché il PMA a 1 parametro crolla sotto DM) e
stabilisce la regola poi ereditata, senza bisogno di rideriverla, sia
dall'anello sia dalla catena.

\textbf{Verifiche.} Eseguito per intero in questa sessione, zero errori.

\textbf{Conclusione.} È il notebook di indagine mirata originale del progetto — il
modello diretto da cui sono stati poi costruiti
\texttt{analisi\_espressivita\_PMA\_anello.ipynb} e l'analogo, non ancora prodotto,
per la catena.

\section*{\texttt{vqe\_ground\_state\_test2.ipynb} + \texttt{ground\_state\_test2.npz}}

VQE al punto di lavoro "test 2" ($b/J=0.35$, $D/J=0.80$): calcola il ground state
con il PMA e ne salva lo stato completo (non solo l'energia) in un file
\texttt{.npz}, per essere riusato come stato iniziale in fasi successive (es.\ le
correlazioni dinamiche).

\textbf{Verifiche.} Eseguito per intero in questa sessione, zero errori. Il file
\texttt{.npz} allegato contiene $E_\text{vqe}=-3.5732114465461704$ contro
$E_\text{esatto}=-3.5732114466175835$ (fidelity $0.9999999999849056$); la
riesecuzione in questa sessione ha rigenerato lo stesso file con
$E_\text{vqe}=-3.5732114465400517$ (fidelity $0.9999999999835552$) — differenze
nell'undicesima cifra, dovute a un run di ottimizzazione diverso ma dello stesso
ordine di precisione.

\textbf{Considerazioni.} \texttt{ground\_state\_test2.npz} non è un file
indipendente: è l'output salvato da questo notebook (contiene \texttt{b}, \texttt{J},
\texttt{D}, \texttt{E0\_exact}, \texttt{E\_vqe}, \texttt{fidelity},
\texttt{psi0\_exact}, \texttt{psi0\_vqe}, \texttt{vqe\_params}) — i due vanno
sempre tenuti insieme.

\textbf{Conclusione.} Fornisce lo stato di riferimento già usato per le correlazioni
dinamiche del dimero.

\section*{Conclusione generale}

La fase VQE del dimero è completa e verificata: sei file (due dei quali, il
notebook di test2 e il suo \texttt{.npz}, sempre accoppiati), tutti coerenti fra
loro e con la documentazione già esistente. Nessun file risulta ridondante o
eliminabile — ciascuno copre un aspetto distinto (produzione, guida, confronto
sistematico, indagine mirata, punto di lavoro specifico). L'unico punto lasciato
aperto è la riesecuzione end-to-end di \texttt{confronto\_ansatz\_entangler.ipynb},
non completata in questa sessione per un problema tecnico, non per un dubbio sul
contenuto.

\end{document}
